\vspace{0.5em}\noindent\textbf{Tự đánh giá}\par

Trong quá trình thực hiện project \textbf{Internship Application
Platform}, mình có cơ hội áp dụng kiến thức về full-stack development,
Cloud architecture, AI integration, realtime communication,
containerization, Kubernetes, CI/CD và observability. Phần tự đánh giá
này tổng hợp mức độ hoàn thành, điểm mạnh, hạn chế và định hướng cải
thiện của mình sau kỳ thực tập.

\vspace{0.5em}\noindent\textbf{Tiêu chí đánh giá}\par

{
\begingroup
\small
\setlength{\tabcolsep}{3pt}
\renewcommand{\arraystretch}{1.15}
\begin{longtable}{@{}
  >{\raggedright\arraybackslash}p{(\linewidth - 10\tabcolsep) * \real{0.1667}}
  >{\raggedright\arraybackslash}p{(\linewidth - 10\tabcolsep) * \real{0.1667}}
  >{\raggedright\arraybackslash}p{(\linewidth - 10\tabcolsep) * \real{0.1667}}
  >{\raggedright\arraybackslash}p{(\linewidth - 10\tabcolsep) * \real{0.1667}}
  >{\raggedright\arraybackslash}p{(\linewidth - 10\tabcolsep) * \real{0.1667}}
  >{\raggedright\arraybackslash}p{(\linewidth - 10\tabcolsep) * \real{0.1667}}@{}}
\toprule\noalign{}
\begin{minipage}[b]{\linewidth}\raggedright
STT
\end{minipage} & \begin{minipage}[b]{\linewidth}\raggedright
Tiêu chí
\end{minipage} & \begin{minipage}[b]{\linewidth}\raggedright
Mô tả
\end{minipage} & \begin{minipage}[b]{\linewidth}\raggedright
Tốt
\end{minipage} & \begin{minipage}[b]{\linewidth}\raggedright
Khá
\end{minipage} & \begin{minipage}[b]{\linewidth}\raggedright
Trung bình
\end{minipage} \\
\midrule\noalign{}
\endhead
\bottomrule\noalign{}
\endlastfoot
1 & \textbf{Kiến thức và kỹ năng chuyên môn} & Hiểu lĩnh vực chuyên môn,
áp dụng kiến thức vào thực tế, sử dụng công cụ và đảm bảo chất lượng
công việc & \checkmark & \(\square\) & \(\square\) \\
2 & \textbf{Khả năng học hỏi} & Khả năng tiếp thu kiến thức mới và học
nhanh & \checkmark & \(\square\) & \(\square\) \\
3 & \textbf{Tính chủ động} & Chủ động nhận việc, tìm hiểu vấn đề và đề
xuất hướng xử lý & \checkmark & \(\square\) & \(\square\) \\
4 & \textbf{Tinh thần trách nhiệm} & Hoàn thành công việc đúng thời gian
và đảm bảo chất lượng & \checkmark & \(\square\) & \(\square\) \\
5 & \textbf{Kỷ luật} & Tuân thủ lịch trình, quy định và quy trình làm
việc & \(\square\) & \checkmark & \(\square\) \\
6 & \textbf{Tinh thần cầu tiến} & Sẵn sàng tiếp nhận phản hồi và cải
thiện bản thân & \checkmark & \(\square\) & \(\square\) \\
7 & \textbf{Giao tiếp} & Trình bày ý tưởng và báo cáo công việc rõ ràng
& \(\square\) & \checkmark & \(\square\) \\
8 & \textbf{Làm việc nhóm} & Phối hợp hiệu quả với đồng đội và tham gia
vào hoạt động nhóm & \checkmark & \(\square\) & \(\square\) \\
9 & \textbf{Tác phong chuyên nghiệp} & Tôn trọng đồng đội, đối tác và
môi trường làm việc & \checkmark & \(\square\) & \(\square\) \\
10 & \textbf{Kỹ năng giải quyết vấn đề} & Xác định vấn đề, đề xuất giải
pháp và thể hiện tư duy sáng tạo & \checkmark & \(\square\) &
\(\square\) \\
11 & \textbf{Đóng góp cho project/team} & Hiệu quả công việc, đóng góp
kỹ thuật và sự ghi nhận từ nhóm & \checkmark & \(\square\) &
\(\square\) \\
12 & \textbf{Tổng quan} & Đánh giá chung cho toàn bộ kỳ thực tập &
\checkmark & \(\square\) & \(\square\) \\
\end{longtable}
\endgroup
}

\vspace{0.5em}\noindent\textbf{Điểm mạnh}\par

\begin{itemize}
\tightlist
\item
  Mình có khả năng tiếp cận vấn đề theo hướng hệ thống, không chỉ tập
  trung vào một chức năng riêng lẻ mà còn xem xét authentication,
  storage, realtime chat, AI processing, deployment và observability.
\item
  Mình chủ động học thêm các dịch vụ AWS và công cụ Cloud-native để hiểu
  cách một ứng dụng có thể vận hành trong môi trường production.
\item
  Mình biết cách liên hệ kiến thức đã học với project thực tế, ví dụ như
  áp dụng S3 presigned URL cho bảo mật tài liệu, Redis adapter cho
  realtime chat, và worker cho các tác vụ xử lý lâu.
\item
  Mình cải thiện kỹ năng viết báo cáo kỹ thuật, trình bày nội dung song
  ngữ và chuẩn hóa cấu trúc tài liệu theo template FCAJ.
\item
  Mình có tinh thần tiếp thu phản hồi và sẵn sàng chỉnh sửa nội dung khi
  phát hiện điểm chưa hợp lý.
\end{itemize}

\vspace{0.5em}\noindent\textbf{Điểm cần cải
thiện}\par

\begin{itemize}
\tightlist
\item
  Mình cần cải thiện khả năng ước lượng thời gian cho từng phần công
  việc, đặc biệt với các phần có nhiều thành phần liên quan như
  Kubernetes, observability và CI/CD.
\item
  Mình cần thực hành sâu hơn về triển khai production trên AWS, bao gồm
  networking, IAM, cost control và monitoring thực tế.
\item
  Mình cần nâng cao kỹ năng viết automated tests để bao phủ tốt hơn các
  luồng auth, application, document, chat và AI.
\item
  Mình cần cải thiện khả năng thiết kế UI/UX để giao diện dễ sử dụng hơn
  cho cả Candidate và HR.
\item
  Mình cần tiếp tục rèn luyện cách giải thích kết quả AI để tránh sử
  dụng điểm matching như một quyết định tuyển dụng duy nhất.
\end{itemize}

\vspace{0.5em}\noindent\textbf{Kế hoạch học tập tiếp
theo}\par

\begin{itemize}
\tightlist
\item
  Ôn tập và thực hành sâu hơn về AWS Well-Architected Framework, đặc
  biệt là security, reliability và cost optimization.
\item
  Tiếp tục học Kubernetes, Helm, Ingress, HPA, PDB, probes và deployment
  strategies.
\item
  Tìm hiểu thêm về observability với Prometheus, Grafana, Loki, Tempo và
  OpenTelemetry.
\item
  Cải thiện kỹ năng backend với FastAPI, SQLAlchemy, Alembic, PostgreSQL
  transaction và security best practices.
\item
  Thực hành CI/CD với GitHub Actions, OIDC, container scanning, secret
  scanning và deployment automation.
\item
  Nghiên cứu thêm về AI application design, prompt engineering, schema
  validation, fallback strategy và human-in-the-loop review.
\end{itemize}

\vspace{0.5em}\noindent\textbf{Định hướng nghề
nghiệp}\par

Sau kỳ thực tập, mình có định hướng tiếp tục phát triển theo hướng
\textbf{Cloud/Backend Engineering} kết hợp với \textbf{AI integration}.
Project này giúp mình nhận ra rằng một hệ thống thực tế không chỉ cần
chạy được chức năng chính, mà còn cần bảo mật, khả năng mở rộng,
observability, tự động hóa triển khai và kiểm soát chi phí.

Trong thời gian tới, mình muốn tiếp tục xây dựng các project có tính
thực tế cao hơn, triển khai trên AWS một cách có kiểm soát và rèn luyện
tư duy thiết kế hệ thống để có thể tham gia tốt hơn vào các dự án
Cloud-native trong môi trường doanh nghiệp.
